\paragraph{Conjugate-coordinate class: exact Landau (nilpotent) resummation of the leading symbol in the background flux, and its effect on the mass-supertrace and order-\(F^4\) terms (conditional; bounded no-go).}
Keep the setting of the preceding paragraphs: Poisson limit, leaf \(\mathbb{CP}^{1,2}\), \(T^a=M^{a4}\), \(H_X=-\Box_X\), the supplied content with no bilinear fermion mass, \(M_{\mathrm{str}}=4\mu^2+n_{\mathrm{ex}}\mu_i^2\), and the nilpotent (leading) symbol, in which the horizontal momenta \(\omega_a=\xi(W_a)\) obey \([\omega_a,\omega_b]=\mathrm iF_{ab}(k)\) with \(F_{ab}\) the constant flux at the centre variable \(k\) of the kernel plane of \(\sigma(\Box)\); \(\eta F\) has eigenvalues \(\pm f_e\) and \(\pm\mathrm if_m\). At four sampled orbit points the ratio \(f_e/f_m\) on the kernel plane equals \(\ell/|y^4|\) in every direction of \(k\) to five digits, which is the closed form of the kernel-plane flux above. Write \(a=f_e\), \(b=f_m\), \(x=sa\), \(y=sb\).

\emph{Orbital factor.} For \(H_0=\eta^{ab}\omega_a\omega_b\) the Gaussian (KMS) two-point function \(G=-\mathrm iFM^{T}(1-M^{T})^{-1}\), \(M=\mathrm e^{2sF\eta}\), gives, per unit horizontal volume, \(\Theta(s)=\operatorname{Tr}\mathrm e^{-\mathrm isH_0}=-\mathrm i\,ab\,\bigl(16\pi^2\sinh x\,\sin y\bigr)^{-1}\), with \(s\to s-\mathrm i0\) (the poles of \(1/\sin y\) lie above the contour, as for the Landau-level sum \(\tfrac{b}{2\pi}\sum_n\mathrm e^{-\mathrm isb(2n+1)}\)). The identity \(\mathrm d\ln\Theta/\mathrm ds=-\mathrm i\eta^{ab}G_{ab}\) holds to \(10^{-10}\) at complex \(s\), and \(\Theta\to-\mathrm i/(16\pi^2s^2)\) for \(s\to0\). The electric pair gives \(1/\sinh\) and the magnetic pair \(1/\sin\).

\emph{Mass-supertrace term.} The term \(+M_{\mathrm{str}}\operatorname{Tr}[(H_X-\mathrm i0)^{-1}]=M_{\mathrm{str}}\,\mathrm i\!\int_0^\infty\!\mathrm ds\,\operatorname{Tr}\mathrm e^{-\mathrm isH_X}\) becomes, per unit horizontal volume and per centre point, \(M_{\mathrm{str}}(16\pi^2)^{-1}(s_0^{-1}+g)\) with
\[
g(a,b)=\int_0^\infty\!\mathrm ds\Bigl[\frac{ab}{\sinh(as)\,\sin(bs-\mathrm i0)}-\frac1{s^2}\Bigr]=\pi\bigl[\,b\,S(b/a)+\mathrm i\,a\,S(a/b)\bigr],\qquad S(r)=\sum_{n\ge1}\frac{(-1)^n}{\sinh(n\pi r)}<0 .
\]
The integrand is analytic in the open fourth quadrant and bounded by \(\exp[-|s|(a\cos\vartheta+b\sin\vartheta)]\) on the ray \(s=|s|\mathrm e^{-\mathrm i\vartheta}\), so \(g\) is finite for \(a,b>0\) and independent of \(\vartheta\); the closed form follows from the two residue sums (poles of \(1/\sinh\) at \(s=-\mathrm in\pi/a\), poles of \(1/\sin\) at \(s=n\pi/b\)) and was confirmed against the ray integral to \(10^{-24}\) at ten pairs \((a,b)\). \(g\) is homogeneous of degree one, \(\operatorname{Re}g<0\) and \(\operatorname{Im}g<0\) for all \(a,b>0\), \(\operatorname{Re}g\to-a\ln2\) for \(b\ll a\) and \(\operatorname{Im}g\to-b\ln2\) for \(a\ll b\). Thus the linear divergence in \(s\) of the Maxwell tail \(c_0\) of the preceding paragraph is cut off at \(s\sim1/f\), and the flux supplies an exponentially convergent infrared end of the proper-time integral.

\emph{Completeness of the \(M_{\mathrm{str}}\) term.} With the flux kept to all orders, the mass-linear term is not \(M_{\mathrm{str}}\operatorname{Tr}[(H_X-\mathrm i0)^{-1}]\) alone: the vector mass multiplies the Zeeman-dressed resolvent, and the complete term is \(\mathrm i\!\int\!\mathrm ds\,\Theta(s)\bigl[M_{\mathrm{str}}+\mu^2(\chi_V-4)\bigr]\), \(\chi_V=2\cosh2x+2\cos2y\). The second part is the all-order version of the constant \(\sum_iw_im_i\operatorname{tr}\Sigma_i\Sigma_i\) listed above and has the growth of the next paragraph.

\emph{Order-\(F^4\) character.} With \(\chi_{\mathrm{str}}=4(\cos y-\cosh x)^2\) (the species sum \(2\cosh2x+2\cos2y+6-2-8\cosh x\cos y\), identity checked to \(10^{-40}\)) the resummed integrand is \(\Theta\chi_{\mathrm{str}}/s\). On the ray of angle \(\vartheta\) one has \(\bigl|\Theta\chi_{\mathrm{str}}\bigr|\sim\exp\bigl[|s|\,|a\cos\vartheta-b\sin\vartheta|\bigr]\) (measured rates agree to \(4\cdot10^{-3}\) at four angles away from \(\vartheta_*\)), and \(\Theta\chi_{\mathrm{str}}\to(ab/\pi^2)\sinh^2[s(a-\mathrm ib)/2]\) at large \(|s|\) (ratio \(1\) to \(10^{-8}\)). The exponential growth therefore occurs on the real axis (the \(\mathrm i\epsilon\) contour) for every \(a>0\), the same exponent holds for \(\Theta(\chi_V-4)\), the exponent is unchanged by the phases \(\mathrm e^{-\mathrm ism_i}\) of real masses, and it vanishes only on the single ray \(\vartheta_*=\arctan(a/b)\), which depends on the flux; on the leaf \(\vartheta_*=\arctan(\ell/|y^4|)\). On that ray the integrand is bounded, but it tends to the constant \(-ab/(2\pi^2)\) plus oscillations, so \(\int\mathrm ds/s\) grows logarithmically in the real part (numerically \(-0.032\) per doubling of the upper limit at \(a=0.7\), \(b=1.3\); the imaginary part converges). No ray gives a convergent integral, and the order-\(F^4\) term has no finite value from the proper-time representation of the class; its non-analytic dependence on \(f_e\) is that of a spectral sum over the complex eigenvalues \(\lambda\pm2\mathrm if_e\) of the vector sector and is not computed.

\emph{Dependence on \(A\).} The leading nilpotent coefficient depends on \(A\) only through the centre-valued flux, so its second variation is \(\tfrac12D^2g\,[Q,Q]\) with \(Q_{ab}=\zeta_Y(k)(\kappa_a\alpha_b-\kappa_b\alpha_a)\); it contains no second derivative of the fluctuation and no term of Einstein--Hilbert form. At four orbit points, averaging over the direction of \(k\) with the window of \(\Delta_{3/4}\), this form is gauge invariant (\(\alpha\propto\kappa\) in its kernel) but is not proportional to \((\kappa\cdot\kappa)\eta^{ab}-\kappa^a\kappa^b\) (relative residual \(0.4\)--\(1.0\) in real and imaginary parts at the points with \(0.3\le f_e/f_m\le0.8\)), and its eigenvalues on the quotient by \(\kappa\) have both signs. The infrared-finite replacement of the Maxwell tail is therefore neither covariant nor sign-definite.

\emph{Einstein--Hilbert sector in a four-dimensional Landau model.} Let \(\Omega_a=\omega_a+\epsilon\sin\phi\,(m_a{}^c\omega_c+n_a)\), \(\phi=\kappa\cdot x\), \(m_a{}^c\kappa_c=0\), on the algebra above with \(F\) constant; the choice \(m_a{}^c=-\alpha_ay^c\), \(n_a=-\alpha_a c\) reproduces the sector of the preceding paragraph with the commutators \([\omega_a,\omega_b]=\mathrm iF_{ab}\) kept, but without the commutator of \(Y\) with \(W_a\) that the Jacobi identity fixes on the leaf, so it is a model and not the leaf. The second order in \(\epsilon\) of \(\operatorname{Tr}\mathrm e^{-\mathrm isH}\) was computed exactly by Wick's theorem in the state \(\mathrm e^{-\mathrm isH_0}\). It reproduces the Maxwell coefficient at \(F=0\) (ratio \(1\) to \(10^{-10}\)), the form \(A\) of the preceding paragraph at \(F=0\) (\(c_1/A=\mathrm i\) to \(10^{-8}\) in five random draws, including the dilaton term), gauge invariance in \(n\) at \(F\ne0\) (\(2\cdot10^{-15}\)), and the constant-field limits (Hessian of \(\Theta\) in \(F\), and \(|\det M|^{-1}\Theta(MFM^T)\) for a constant vielbein). At \(F\ne0\) the coefficient \(c=A/16\pi^2\) of \(\mathrm ds/s\) is unchanged, the integrand decays like \(\mathrm e^{-|s|(a\cos\vartheta+b\sin\vartheta)}\), and the finite part \(\mathrm{Fin}=\lim[\int_{s_0}^\infty T_2\,\mathrm ds-c\ln(1/s_0)]\), in units of \(c\), is about \(7.6\), \(8.1\), \(6.8\) in three Lorentz frames of \(F\) at \((a,b)=(0.7,1.3)\) (data \(\kappa,\alpha,y\) rotated, \(A\) fixed), and \(7.6\), \(5.9\), \(7.3\) in three frames at \((a,b)=(1.92,1.0)\). The flux therefore turns the infrared end of \(\int\mathrm ds/s\) into a finite remainder of the size of the logarithm that depends on the orientation of \(F\) relative to \(\kappa,\alpha,y\); the Maxwell part of the same sector changes by up to \(40\%\) between frames. The coefficient of the logarithm is the flux-independent \(A\) of the preceding paragraph.

This is a statement for the Poisson limit, one loop, the nilpotent (leading) symbol resummed in a constant flux at the centre variable, the Lorentzian \(\mathrm i\epsilon\) proper time, the supplied content, and numerical evaluation at sampled points; the closed form of \(g\) is checked numerically, not proved, and the Landau model for the Einstein--Hilbert sector is not the leaf. It does not include the next-to-leading symbol of the hypoelliptic operator with the commutators of \(Y\), the dependence of the window on the flux, non-perturbative masses, or the spectral definition of the vector sector, and it identifies no induced form with the emergent metric of Chapter~\ref{ch:vierbein-emergence}, makes no Wick-rotation claim and gives no value or sign of \(G_{\mathrm{eff}}\). This changes no ledger label; Problem~F remains open.
